\section{Models and Evaluation Harness}
\label{sec:models}

\textbf{Roster and eligibility.} Eight models are scored. A model may forecast a
row only if its exact published version was released before the row's outcome
could be known; in practice every model was released before its first window
opened. The four models published between August 2025 and January 2026
run on all three windows. The four published between 7 and 24 April 2026 are ineligible for \winA (February 2026) and eligible for \winC (May to June) and \winB (July to August). The rule is a
property of dates rather than of a leakage test that might fail silently. The
full roster is in Appendix~\ref{app:roster}.

\textbf{Hosts and settings.} \mGptFiveTwo and \mGptFiveFive are served by their
developer; OpenAI returned dated snapshots on every call, confirming the
checkpoint did not change. Every open-weight model is pinned to a single host
with fallback routing off, and the served model string is checked on every reply.
Settings are fixed: temperature 0 where the host accepts it, one sample per row, 4{,}000-token reply limit.
Thinking is a reasoning-effort level of \emph{medium} on the OpenAI models and
\mGptOss, and a 1{,}000-token budget on the others; the asymmetry cannot be
removed (Section~\ref{sec:limitations}).

\textbf{Prompt.} Every model receives the same template
(Appendix~\ref{app:prompt}), requiring a final \code{PROBABILITY:} line. It
contains no price, bid, ask, depth or book. All scored calls use this one template.

\textbf{Forced answers.} An unparseable reply is re-asked once with a stricter
instruction; a second failure is recorded as failed. Across all scored runs 79
rows failed, all from four models, 42 inside the primary set. On the test split, filling every failure with 0.5 moves no Brier edge by more
than 0.0011; filling each with the worst possible answer moves none by more than
0.0077, always in the direction
that makes the models look worse. The re-ask was needed on 4 to 14 per cent of rows for \mKimiKTwoFive, about
5 per cent for \mDeepSeekVFourPro and at most 3 per cent for every other model.

\textbf{Leak checks at call time.} Two checks run before every call. The first
rejects a row whose news capture time is not strictly before its forecast moment.
The second scans for a dateline two or more days after the forecast moment,
besides the benchmark's own flag, which is anchored on the capture time
(Appendix~\ref{app:newsattach}). A row that would leak is never sent.

\textbf{Memory probe.} Every news-bearing row is asked again with the news
removed. The pair measures what the news was worth and flags rows a model may
have memorised; flag rates over whole probe runs are 0.8 to 2.3 per cent. Appendix~\ref{app:memory}
gives the rule and a validation exercise whose result is genuinely inconclusive
(Section~\ref{sec:limitations}).

\textbf{Scale and cost.} The study is 97{,}465 model calls, costing about \$620.
